\noindent\textbf{Résumé.}
\par\smallskip
La sélection conjointe exceptionnelle et temporelle restreint le domaine de \(19\,446\) candidats à \(79\) par incidence et à un unique \(K\) par compatibilité temporelle. Le répertoire, les orientations et les filtres qui produisent cette unicité sont conservés.


On détermine la constante holographique de couplage arithmético--géométrique au moyen d’un registre ordonné \(K\), de son évaluation rationnelle et de ses représentations incidentielles, et on démontre une loi de conservation qui permet de retrouver l’information à travers ses transformations. La base mille et le système de coordonnées de phase étant fixés, la lecture \(\kappa=N_K/(1000^{12}-1)\) conserve exactement les douze composantes du registre. Son orbite cyclique forme une base et fournit une identification stable des opérateurs linéaires : on utilise les réponses à la base complète ou une réponse vectorielle lorsqu’il existe une covariance cyclique. Ainsi, la structure qui détermine la valeur participe également à la reconstruction des opérations qui agissent sur elle.

La construction part d’un noyau formel appelé holographie modulaire triadique, ci-après HMT. Il comprend \emph{All Possible Paths} (APP), structure arithmétique de chemins sur un support toroïdal avec un graphe de Cayley et des évaluations additive et multiplicative distinguées ; le TRIT, signature orientée de valeurs \(\{-1,0,+1\}\) ; et le \emph{Topological Phase Kernel} (TPK), qui compose sélection, transport et mémoire. Le prolongement nonadique conserve l’état enrichi dont proviennent les coordonnées arithmétiques et les incidences de frontière. Le registre \(K\) intervient dans la clôture de \(\alpha\) avec les coordonnées HMT de \(\pi,\varphi,e\) ; sa réalisation incidencielle transporte la même information ordonnée vers les espaces spécifiés.

Le résultat général de conservation est une identité bilatérale pour deux isométries avec domaine et codomaine communs, valable en dimension arbitraire. Le raffinement entier avec retenue possède un inverse exact sur les étiquettes admissibles et un relèvement unitaire qui se compose avec cette identité. Dans la spécialisation nonadique, les poids \(72/73\) et \(1/73\) déterminent la transformation des observables ; les invariants sont exactement ceux qui commutent avec le transport relatif. Une politique explicite d'archivage transfère toute la norme à la mémoire infinie et offre une reconstruction par un adjoint de norme un, ainsi que des inverses finis avec des bornes uniformes en profondeur. Deux observations complémentaires récupèrent avec une borne optimale le registre centré, à partir duquel est reconstruit le projecteur dimensionnel spécifié. La conservation s'étend à chaque degré extérieur admissible, y compris aux contributions mixtes de lecture et de mémoire.

L’extension réversible relie déformation, holonomie et action ; les réalisations électromagnétiques et gravitationnelles conservent leurs opérateurs et leurs normalisations explicites. Pour le protocole symplectique déclaré, avec deux entrées gaussiennes pures, centrées et identiques et un nombre fini de modes, la composante antilinéaire de la mémoire détermine le spectre symplectique réduit, en maintenant toutes les occupations de Fock. Une boucle elliptique--parabolique donne une pureté \(9/11\), un spectre \((9/10)10^{-j}\) et une entropie adimensionnelle \((10/9)\log10-\log9\). Le résultat conjoint caractérise ce que conserve une transformation, ce que modifie une observation et les données qui permettent de l’inverser.


\medskip
% BEGIN RECIPROCIDAD 20260918
La direction récupérable de \(K\) détermine en outre un flux positif de déterminant un, conserve les strates positroïdales et produit des réseaux duaux de covolume fixe à symétrie ternaire. Sa lecture thêta--Mellin conserve les pôles et les résidus, tandis que la variation spectrale réside dans une différence entière. Le contraste angulaire engendré se réalise comme coordonnée de Klein ; avec l’orientation conservée, le rayon retrouve la section d’action et réexprime la formule électronique complète. Les preuves relient cette géométrie à l’archive bilatérale et fournissent des bornes quantitatives de reconstruction.

% END RECIPROCIDAD 20260918

% BEGIN R27_FRAMING_SUMMARY_ES
La base émettrice de cette récupération est explicitée au moyen de l'alphabet de \(42\) blocs décimaux, des raccordements régionaux et d'un exemple qui sépare mémoire ordonnée et histogramme. La comparaison des jauges détermine la normalisation de l'observable multiplicatif, et les contrôles de clôture séparent orientation, condition terminale et position du registre. À la frontière conjointe, saturation et neutralité duale laissent deux représentants ; l'orientation fixée détermine la section sélectionnée. Ces résultats identifient les données que reçoivent ensuite les lois de conservation et de récupération.
% END R27_FRAMING_SUMMARY_ES

\noindent\textbf{Mots-clés :} holographie modulaire triadique ; mécanique dimensionnelle ; extrême et moyenne raison holographique ; couplage arithmético--géométrique ; raffinement réversible ; mémoire opératoire ; holonomie ; incidence exceptionnelle ; conservation extérieure ; transport bilatéral ; observables invariantes ; récupération stable ; action hamiltonienne ; covariance gaussienne ; entropie d’intrication.


\clearpage
\noindent\textbf{Abstract.}\par\smallskip
Joint exceptional and temporal selection restricts the domain of \(19\,446\) candidates to \(79\) through incidence and to a unique \(K\) through temporal compatibility. The repertoire, orientations and filters producing this uniqueness are retained explicitly.

The holographic arithmetic--geometric coupling constant is determined through an ordered register \(K\), its rational evaluation and its incidence representations, and a conservation law is proved that permits information recovery through these transformations. With base one thousand and the phase chart fixed, the reading \(\kappa=N_K/(1000^{12}-1)\) preserves all twelve register components exactly. Its cyclic orbit forms a basis and provides stable identification of linear operators: the data are the responses to the complete basis, or one vector response under cyclic covariance. The structure that determines the value thus also participates in reconstructing the operations acting on it.

The construction starts from a formal kernel called Triadic Modular Holography, hereafter HMT. It comprises \emph{All Possible Paths} (APP), an arithmetic path structure on a toroidal support with a Cayley graph and distinct additive and multiplicative evaluations; TRIT, an oriented signature taking values in \(\{-1,0,+1\}\); and the \emph{Topological Phase Kernel} (TPK), which composes selection, transport and memory. Nonadic prolongation preserves the enriched state from which arithmetic coordinates and boundary incidences arise. The register \(K\) enters the closure of \(\alpha\) together with the HMT coordinates of \(\pi,\varphi,e\); its incidence realization transports the same ordered information to the specified spaces.

The general conservation result is a bilateral identity for two isometries with a common domain and codomain, valid in arbitrary dimension. Integer refinement with carry has an exact inverse on admissible labels and a unitary lift that composes with this identity. In the nonadic specialization, the weights \(72/73\) and \(1/73\) determine the transformation of observables; the invariants are precisely those commuting with the relative transport. An explicit archive policy transfers the entire norm to infinite memory and yields reconstruction by an adjoint of norm one, together with finite inverses whose bounds are uniform in depth. Two complementary observations recover the centred register with an optimal bound, from which the specified dimensional projector is reconstructed. Conservation extends to every admissible exterior degree, including mixed readout--memory contributions.

The reversible extension relates deformation, holonomy, and action; the electromagnetic and gravitational realizations retain their explicit operators and normalizations. For the declared symplectic protocol with two identical centred pure Gaussian inputs and finitely many modes, the antilinear component of memory determines the reduced symplectic spectrum, retaining all Fock occupations. An elliptic--parabolic loop gives purity \(9/11\), spectrum \((9/10)10^{-j}\), and dimensionless entropy \((10/9)\log10-\log9\). Taken together, the results characterize what a transformation preserves, what an observation changes, and which data permit its inversion.

\medskip
% BEGIN RECIPROCIDAD 20260918
The recoverable direction of \(K\) also determines a positive determinant-one flow, preserves positroid strata, and produces dual lattices of fixed covolume with ternary symmetry. Their theta--Mellin readout preserves poles and residues, while spectral variation lies in an entire difference. The generated angular contrast is realized as a Klein coordinate; with orientation retained, the radius recovers the action section and re-expresses the complete electronic formula. The proofs connect this geometry with the bilateral archive and provide quantitative reconstruction bounds.

% END RECIPROCIDAD 20260918

% BEGIN R27_FRAMING_SUMMARY_EN
The emitting basis of this recovery is made explicit through the alphabet
of \(42\) decimal blocks, regional junctions, and an example separating
ordered memory from its histogram. Gauge comparison determines the
normalization of the multiplicative observable, while closure controls
separate orientation, terminal condition, and register position. At the
joint boundary, saturation and dual neutrality leave two representatives;
the fixed orientation determines the selected section. These results
identify the data subsequently received by the conservation and
recovery laws.
% END R27_FRAMING_SUMMARY_EN

\noindent\textbf{Keywords:} Triadic Modular Holography; Dimensional Mechanics; holographic extreme and mean ratio; arithmetic--geometric coupling; reversible refinement; operator memory; holonomy; exceptional incidence; exterior conservation; bilateral transport; invariant observables; stable recovery; Hamiltonian action; Gaussian covariance; entanglement entropy.
